\chapter{Gefilterte Struktur-Topologieoptimierung}\label{theory:topology}
Die stationäre Tet4-Mechanik löst für feste Totlasten \(K(\rho)u_l=f_l\).
Homogene Lagerungen sind Voraussetzung der selbstadjungierten
Nachgiebigkeit \(C=\sum_l w_l f_l^T u_l=2\sum_l w_l U_l\) in Joule.
Das physische Volumen ist \(v=\sum_e V_e\rho_e/\sum_e V_e\).
Ein positiver Gewichtsfaktor wird nicht automatisch normiert.
Für Zellzentren \(c_e\) und Radius \(r>0\) in Metern gilt
\[
 H_{ej}=\max(0,r-\lVert c_e-c_j\rVert)V_j,\qquad
 P_{ej}=H_{ej}/\sum_k H_{ek},\qquad \rho=Px.
\]
Feste Entwurfswerte sind 0 oder 1; nach der Filterung werden ihre physischen
Zeilen exakt überschrieben. Mit der freien Diagonalmaske \(M\) ist die
Jacobi-Matrix daher \(J=MPM\). Das berücksichtigt Randzeilen,
ungleiche Zellvolumina und konstante Beiträge fester Nachbarzellen.
SIMP setzt \(K_e=a_e K_e^0\),
\(a_e=\epsilon+(1-\epsilon)\rho_e^p\), standardmäßig \(p=3\),
\(\epsilon=10^{-3}\). Auch eine Leerzelle hat numerische Reststeifigkeit.
Mit der gewichteten Zellenergie \(U_e=\sum_l w_l U_{el}\) liefert
Differenzieren des Zustands bei festen Kräften
\[
 \frac{\partial C}{\partial\rho_e}
 =-\frac{2U_e}{a_e}p(1-\epsilon)\rho_e^{p-1},\qquad
 \nabla_x C=J^T\nabla_\rho C,\quad
 \nabla_x v=J^TV/\sum_e V_e.
\]
SLSQP minimiert \(C/C_0\) mit Bounds und linearer Volumenrestriktion.
Die Koordinaten \(z=x/\sqrt{n_{\rm frei}}\) konditionieren die Schrittnorm;
Gradienten und Schranken werden vollständig transformiert. Ein lokaler Stopp
beweist keine globale Optimalität. Der lineare Filter ist eine Regularisierung,
keine garantierte Mindestwandstärke. Zwischenmaterial und Netzeinfluss müssen
separat ausgewiesen werden.

Für eine diskrete Geometrie werden Zellen über drei feste Schwellen gewählt.
Getrennte Teile sind ein No-go. Ein ausdrücklich optionaler zweiter Weg ergänzt
Flächenpfade mit Kosten \(V(1-\rho+10^{-3})\) im Ursprungsnetz, ohne feste
Leergebiete aufzufüllen. Alle Ergänzungen werden protokolliert. Danach rechnet
der unveränderte Mechaniksolver Vollmaterial, dieselben Lasten und Lagerungen.
Geometrievolumen und Nachgiebigkeit erhalten eigene Grenzen; sie sind kein
Nachweis der ursprünglichen Dichtevolumenrestriktion oder der Festigkeit.
